% =====================================================
% 题型：解三角形恒等式与三角形判定多选题 (q_triangle_identities_three.tex)
% =====================================================
% 难度：★★ (中档)
% =====================================================
\newcommand{\qTriangleIdentitiesThreeStem}{在 \(\triangle ABC\) 中，若 \(\sin A=2\sin B\cos C\)，则下列结论正确的是 \hfill （\quad）}

\newcommand{\qTriangleIdentitiesThreeOptions}{%
  \mcoptions[1]{\(A=2B\)}{\(a=2b\cos C\)}{\(a^2=b^2+c^2\)}{\(\triangle ABC\) 为直角三角形}%
}

\newcommand{\qTriangleIdentitiesThreeAnswer}{B}

\newcommand{\qTriangleIdentitiesThreeSolution}{%
  由正弦定理 \(\dfrac{a}{\sin A}=\dfrac{b}{\sin B}\)，有
  \[
    \frac{a}{b}=\frac{\sin A}{\sin B}=2\cos C.
  \]
  所以
  \[
    a=2b\cos C,
  \]
  故 B 正确。\\
  进一步代入余弦定理
  \[
    c^2=a^2+b^2-2ab\cos C,
  \]
  由 \(a=2b\cos C\) 得
  \[
    c^2=b^2,
  \]
  因此 \(b=c\)，进而 \(B=C\)。结合三角形内角和，可继续分析角关系，但不能仅由原条件直接写出 \(A=2B\) 或断言为直角三角形。%
}

\newcommand{\qTriangleIdentitiesThreeDiagnosis}{%
  看到三角函数条件时，优先用正弦定理把角函数比例翻译成边长关系；不要直接把 \(\sin A=2\sin B\cos C\) 误读为角的倍数关系。%
}

\newcommand{\qTriangleIdentitiesThreeTeachingNotes}{%
  适合训练“角条件 → 边条件 → 再回到几何结构”的双向翻译。第一动作应是识别 \(\sin A/\sin B=a/b\)。%
}
